\documentclass[10pt,a4paper]{amsart}
\usepackage[margin=19mm]{geometry}
\usepackage{amsmath,amssymb,mathtools}
\usepackage[scr=rsfs,cal=euler]{mathalfa}
\usepackage{xfrac}
\usepackage[unicode,hidelinks,bookmarks=false]{hyperref}
\newcommand{\ie}{i.e.\ }
\newcommand{\cf}{cf.\ }
\newcommand{\resp}{respectively\ }
\newcommand{\Subsection}{\subsection}
\providecommand{\nobreakdash}{\nobreak\hbox{-}}
\setlength{\emergencystretch}{2em}
\multlinegap0pt
\usepackage{bm}
\usepackage{bbm}
\usepackage{dsfont}
\usepackage{xspace}
\usepackage{cancel}
\usepackage{mathtools}
\usepackage{amsthm}
\makeatletter
\@ifundefined{proposition}{\newtheorem{proposition}{Proposition}}{}
\@ifundefined{theorem}{\newtheorem{theorem}[proposition]{Theorem}}{}
\makeatother
\DeclareMathOperator{\dist}{dist}
\newcommand{\E}{\mathbb E}
\newcommand{\Ll}{\left}
\newcommand{\Rr}{\right}
\newcommand{\R}{\mathbb{R}}
\newcommand{\al}{\alpha}
\newcommand{\de}{\delta}
\newcommand{\dr}{\partial}
\newcommand{\eps}{\varepsilon}
\newcommand{\ep}{\varepsilon}
\newcommand{\msc}{\mathscr}
\newcommand{\td}{\widetilde}
\renewcommand{\bar}{\overline}
\renewcommand{\d}{\mathrm{d}}
\renewcommand{\ge}{\geqslant}
\renewcommand{\le}{\leqslant}
\title[Oscillating solutions to the mean-field Langevin descent-asc]{Oscillating solutions to the mean-field Langevin descent-ascent flow}
\author{Jean-Christophe Mourrat, Loucas Pillaud-Vivien}
\date{Source: arXiv:2604.11134v1 (CC BY 4.0)}
\begin{document}
\maketitle

\noindent\emph{Source and licence.} Oscillating solutions to the mean-field Langevin descent-ascent flow. Version arXiv:2604.11134v1 (CC BY 4.0), distributed under the Creative Commons Attribution 4.0 International licence, as recorded in the arXiv licence metadata for this exact version and shown on the abstract page https://arxiv.org/abs/2604.11134v1. Full rights notice: licenses/arxiv-2604.11134-CC-BY-4.0.txt.\par\smallskip

\noindent\textit{Editorial scope.} Counted unit: the Theorem whose source label is \texttt{t.bootstrap} in the article (source0.tex lines 521--538, source label \texttt{t.bootstrap}), delivered together with the article's own complete original proof of it (source0.tex lines 539--651, one \texttt{proof} environment of 113 lines). No mathematics is written, repaired or re-derived: statement and proof are the article's own text line for line. Required context retained uncounted: e.particleODE (lines 263--266); e.def.F (lines 268--271); p.square (lines 310--321); p.lyapunov (lines 330--336); e.ev.mn (lines 467--470); p.moment.bounds (lines 477--483); e.ito.pmoment (lines 486--493). Every definition, display and label of those lines is retained unchanged. Omitted from this package and not redistributed: 1--262, 267--267, 272--309, 322--329, 337--466, 471--476, 484--485, 494--520, 652--739. Nothing inside a retained excerpt is omitted or altered: every retained body is delivered exactly as the article's lines are written, so no omission receipt of any kind accompanies this package. Citations: the retained excerpts carry no citation command at all, so no bibliography entry of the article is transcribed and no reference is redistributed. Reference adaptation: none was needed and none was made; every reference inside the delivered unit resolves inside it, so no unresolved reference or citation is produced. Printed slips kept literal: no printed word, symbol, hypothesis or number of the counted statement or of its proof was altered. Layout and class: the article's own class is \texttt{amsart} with packages including fontenc, lmodern, geometry, esint, amssymb, graphicx, subcaption, MnSymbol, mathtools, hyperref, microtype, bm, dsfont, mathrsfs; the wrapper is the lane's pinned \texttt{amsart} editorial wrapper at 10pt on A4, which restates the theorem-like environments the retained text needs and adds the article's own macro definitions that the retained text needs (\texttt{\textbackslash E}, \texttt{\textbackslash Ll}, \texttt{\textbackslash R}, \texttt{\textbackslash Rr}, \texttt{\textbackslash al}, \texttt{\textbackslash bar}, \texttt{\textbackslash d}, \texttt{\textbackslash de}, \texttt{\textbackslash dist}, \texttt{\textbackslash dr}, \texttt{\textbackslash ep}, \texttt{\textbackslash eps}, \texttt{\textbackslash ge}, \texttt{\textbackslash le}, \texttt{\textbackslash msc}, \texttt{\textbackslash td}), with extra wrapper packages (bm, bbm, dsfont, xspace, cancel, mathtools). The delivered numbering is the unit's own. Assets: no retained span contains a figure environment, a graphics-inclusion command, a PDF-page inclusion, a table or a TikZ picture; no image file is copied into this package, the package declares no asset dependency and no image file is redistributed with it. Licence: version arXiv:2604.11134v1 of this article is distributed under the Creative Commons Attribution 4.0 International licence, as recorded in the arXiv licence metadata for this exact version and shown on the abstract page https://arxiv.org/abs/2604.11134v1. A full rights notice is delivered at licenses/arxiv-2604.11134-CC-BY-4.0.txt. Characters: every retained excerpt is pure ASCII, so the delivered engine, XeLaTeX with its default Latin Modern text font, reports no missing character.
\begin{equation}
\label{e.particleODE}
\dr_t m = -\al n + m - \frac{m^3}{3}, \qquad \dr_t n = \al m + n - \frac{n^3}{3}.
\end{equation}

\begin{equation}  
\label{e.def.F}
F_\al (m,n) := \Ll(-\al n+m-\frac{m^3}{3}, \al m + n - \frac{n^3}{3}\Rr).
\end{equation}

\begin{proposition}[square average]
\label{p.square} Consider the system \eqref{e.particleODE} initialized at a point in the cycle $\msc C_\al$, and recall that we denote by $T_\al$ the minimal period of the resulting trajectory. We have
\begin{equation}
\label{e.mean}
\frac 1 {T_\al} \int_0^{T_\al} m_t \, \d t = \frac 1 {T_\al} \int_0^{T_\al} n_t \, \d t = 0
\end{equation}
and
\begin{equation}  
\label{e.square}
\frac 1 {T_\al} \int_0^{T_\al} m_t^2 \, \d t = \frac 1 {T_\al} \int_0^{T_\al} n_t^2 \, \d t \ge \frac 3 2.
\end{equation}
\end{proposition}

\begin{proposition}[Lyapunov function]
\label{p.lyapunov}
Let $g(z) := \dist^2(z,\msc C_\al)$ denote the squared distance between $z \in \R^2$ and the limit cycle $\msc C_\al$. There exist $\delta, c > 0$ such that for every $\al$ sufficiently large and $z \in \R^2$, we have
\begin{equation*}  %\label{e.}
g(z) \le \de  \quad \implies \quad (F_\al \cdot \nabla g)(z) \le -c  g(z).
\end{equation*}
\end{proposition}

\begin{equation}  
\label{e.ev.mn}
\dr_t m_t = -\alpha n_t + m_t - \frac{\E[X_t^3]}{3}, \qquad \dr_t n_t = \alpha m_t + n_t - \frac{\E[Y_t^3]}{3},
\end{equation}

\begin{proposition}[Moment bounds]
\label{p.moment.bounds} 
For every even integer $k \ge 2$ and provided that $\E[X_0^k + Y_0^k]$ is finite, there exists a constant $C_k$ depending only on $k$ and $\E[X_0^k + Y_0^k]$ such that for every $\al \ge 0$, $\ep \in (0,1]$, and $t \ge 0$, we have
\begin{equation*}  %\label{e.}
\E[\td X_t^k + \td Y_t^k] \le C_k.
\end{equation*}
\end{proposition}

\begin{multline}  
\label{e.ito.pmoment}
\frac 1 k \dr_t \E[\td X_t^k] 
 = \E \Ll[\td X_t^{k-1}\Ll( (1-m_t^2) \td X_t -m_t \Ll(\td X_t^2 - \E[\td X_t^2] \Rr)
 - \frac 1 3 \Ll( \td X_t^3 - \E[\td X_t^3] \Rr)\Rr) \Rr] 
 \\ 
 + \ep (k-1) \E[\td X_t^{k-2}].
\end{multline}

\begin{theorem}[Induction]
\label{t.bootstrap}
There exists a constant $C < +\infty$ such that the following holds for every $k  > \frac 3 2$, $\alpha < +\infty$ sufficiently large, $\eps > 0$ sufficiently small, initialization $X_0, Y_0$ with $X_0^2 + Y_0^2 \le 8$, and $t_0 \ge 0$. If
\begin{equation}  
\label{e.bootstrap.ass}
\E[\td X_{t_0}^2]\le \al^{-k}, \quad \E[\td Y_{t_0}^2]\le \al^{-k}, \quad \dist(p_{t_0}, \msc C_\al) \le  \alpha^{1 - k},
\end{equation}
then the same inequalities are valid at time $t_0 + T_\al$, that is,
\begin{equation}  
\label{e.bootstrap.ccl}
\E[\td X_{t_0+T_\al}^2]\le \al^{-k}, \quad \E[\td Y_{t_0+T_\al}^2]\le \al^{-k}, \quad \dist(p_{t_0+T_\al}, \msc C_\al) \le \alpha^{1 - k},
\end{equation}
and moreover, letting $\bar p_{t_0}$ denote the point on $\msc C_\al$ that is closest to $p_{t_0}$ and $(\bar p_t)_{t \ge t_0}$ denote the solution to the deterministic system in \eqref{e.particleODE}, we have for every $t \in [t_0,t_0+T_\al]$ that
\begin{equation}  
\label{e.bootstrap.cycle}
|p_t - \bar p_t| \le C \al^{1 -k}.
\end{equation}
\end{theorem}

\begin{proof}
We fix $k > \frac 3 2$. Under our assumption that the initial condition must satisfy $X_0^2 + Y_0^2 \le 8$, we can bound all moments of $X_t$ and $Y_t$ uniformly over $\alpha \ge 0$, $\ep \in (0,1]$, and $t \ge 0$. In particular, we have $|p_t|^2 = |m_t|^2 + |n_t|^2 \le 8$ for all $t \ge 0$. From now on, we assume that \eqref{e.bootstrap.ass} holds. 

\medskip

\noindent \emph{Step 1.} In this step, we show that for $\alpha$ sufficiently large and $\eps > 0$ sufficiently small, we have for every $t \in [t_0, t_0 + T_\al]$ that $\E[\td X_t^2] \le 2\al^{-k}$ and $\E[\td Y_t^2] \le 2\al^{-k}$. Using that $|m_t|^2 \le 8$ and the It\^o identity \eqref{e.ito.pmoment}, we have that
\begin{equation*}  %\label{e.}
\frac 1 2 \dr_t \E[\td X_t^2] \le  \E[\td X_t^2] + 3 \E[|\td X_t|^3] + \ep.
\end{equation*}
Since by Proposition~\ref{p.moment.bounds} all moments of $\td X_t$ are bounded, we can appeal to H\"older's inequality to find a constant $C < +\infty$ such that $\E[|\td X_t|^3] \le C \E[\td X_t^2]^{1- \frac 1 {2k}}$, so that
\begin{equation}  
\label{e.drX2.crude.bound}
\frac 1 2 \dr_t \E[\td X_t^2] \le  \E[\td X_t^2] + C \E[\td X_t^2]^{1-\frac 1 {2k}} + \ep.
\end{equation}
Denoting $T := \sup \Ll\{ t \in [t_0, t_0 + T_\al] \ : \  \E[\td X_t^2] \le 2 \al^{-k} \Rr\}$, using that $T_\al \le 4\pi/\al$ and redefining $C < +\infty$ as necessary, we have for every $t \in [t_0, T]$ that 
\begin{equation*}  %\label{e.}
\E[\td X_t^2] \le \al^{-k} + C \al^{-1} \Ll( \al^{-k} + \al^{\frac 1 2 - k} + \ep \Rr)  .
\end{equation*}
For $\alpha$ sufficiently large and $\ep > 0$ sufficiently small, we can make sure that the right-hand side above is smaller than $2\al^{-k}$, and therefore that $T = t_0 + T_\alpha$ as desired. The argument for $\E[\td Y_t^2]$ is identical. 

\medskip

\noindent \emph{Step 2.} In this step, we show that provided that $\alpha$ is sufficiently large, we have for every $t \in [t_0, t_0 + T_\al]$ that 
\begin{equation}  
\label{e.dist.unif.time}
\dist(p_t, \msc C_\al) \le \al^{1-k}.
\end{equation}
Recall the notation $F_\al$ in~\eqref{e.def.F}, and the identity \eqref{e.ev.mn} which we can rewrite as
\begin{equation}  
\label{e.evol.p.R}
\dr_t p_t = F_\al(p_t) + R_t, 
\end{equation}
where we have set
\begin{align*}  %\label{e.}
R_t & := \frac 1 3 (m_t^3 - \E[X_t^3], n_t^3 - \E[Y_t^3]) 
\\
& =  \Ll(-m_t \E[\td X_t^2] -\frac 1 3 \E[\td X_t^3], -n_t \E[\td Y_t^2] - \frac 1 3 \E[\td Y_t^3]\Rr).
\end{align*}
Recalling that $\td X_t$ has bounded moments of every order, and using the result of the previous step, we can argue as in the sentence above \eqref{e.drX2.crude.bound} to find a constant $C < +\infty$ such that for every $t \in [t_0, t_0 + T_\al]$,
\begin{equation}  
\label{e.bound.R}
|R_t| \le C \al^{\frac 1 2 - k}.
\end{equation}
Recall that we write $g(z) = \dist^2(z, \msc C_\al)$. We have
\begin{equation*}  %\label{e.}
\dr_t (g(p_t)) = (F_\al \cdot \nabla g)(p_t) + R_t \cdot \nabla g(p_t). 
\end{equation*}
Let $\delta, c > 0$ be as given by Proposition~\ref{p.lyapunov}, and let $T := \sup \{ t \in [t_0, t_0 + T_\al] \ : \ g(p_t) \le \de\}$. For every $t \in [t_0, T]$, we can appeal to this proposition to bound the first term in the previous display by $-c g(p_t)$. 
Using that $|\nabla g(z)| = 2  \sqrt{g(z)}$, the Cauchy-Schwarz and Young's inequalities, we have
\begin{equation*}  %\label{e.}
|R_t \cdot \nabla g(p_t)| \le 2 \sqrt{g(p_t)} |R_t| \le \frac{c}{2} g(p_t) + \frac {2|R_t|^2}{c}.
\end{equation*}
Combining these and enlarging the constant $C < +\infty$, we obtain that for every $t \in [t_0, T]$,
\begin{equation*}  %\label{e.}
\dr_t (g(p_t)) \le -\frac c 2 g(p_t) + C \al^{1-2k}.
\end{equation*}
Integrating this (and updating $C$) yields that for every $t \le T$, 
\begin{equation}
\notag
g(p_t)  \le e^{-\frac{c}{2}t} \al^{2-2k} + C \al^{1-2k} (1 - e^{-\frac c 2 t}).
\end{equation}
By taking $\al \ge C$, we ensure that the right-hand side of the inequality above remains below $\al^{2-2k}$. In particular, since $k > 1$, we have $T = t_0 + T_\al$ and the announced inequality \eqref{e.dist.unif.time} is proved.

\medskip


\noindent \emph{Step 3.} For $\alpha$ sufficiently large, the point in $\msc C_\al$ that is closest to $p_{t_0}$ is well-defined; we denote it by $\bar p_{t_0} = (\bar m_{t_0}, \bar n_{t_0})$, and let $\bar p_t = (\bar m_t, \bar n_t)_{t \ge t_0}$ denote the solution to the system \eqref{e.particleODE} with this initialization. In this step, we show that for a sufficiently large constant $C < +\infty$, the inequality \eqref{e.bootstrap.cycle} is valid. We denote the $L^\infty$ norm of $\nabla F_\al$ on the ball of radius $\sqrt{8}$ by $L_\al$. Recalling that $(p_t)_{t \ge t_0}$ and $(\bar p_t)_{t \ge t_0}$ both take values in this ball, as well as \eqref{e.evol.p.R} and \eqref{e.bound.R}, we have for almost every $t \in [t_0, t_0 + T_\al]$ that
\begin{equation*}  %\label{e.}
\dr_t \Ll(|p_t - \bar p_t|\Rr) \le L_\al |p_t - \bar p_t| + C \al^{\frac 1 2 - k},
\end{equation*}
with
\begin{equation*}  %\label{e.}
|p_0 -\bar p_0| \le \al^{1-k}.
\end{equation*}
Integrating this, we find that for every $t \in [t_0, t_0 + T_\al]$,
\begin{equation*}  %\label{e.}
|p_t -\bar p_t| \le e^{L_\al T_\al} \Ll(\al^{1-k} + C \al^{\frac 1 2 - k}\Rr) .
\end{equation*}
The conclusion follows by noting that $L_\al \le C \al$ while $T_\al \le 4\pi/\al$. 


\medskip

\noindent \emph{Step 4.} We now show that $\E[\td X_{t_0 + T_\al}^2] \le \al^{-k}$ and $\E[\td Y_{t_0 + T_\al}^2] \le \al^{-k}$. As preparation for this, we deduce from the previous step and Proposition~\ref{p.square} that
\begin{equation}  
\label{e.nice.m}
\Ll| \frac 1 {T_\al} \int_{t_0}^{t_0 + T_\al} m_t \, \d t \Rr| \le C \al^{1-k}, \qquad \frac 1 {T_\al} \int_{t_0}^{t_0 + T_\al} m_t^2 \, \d t \ge \frac 5 4.
\end{equation}
Using \eqref{e.ito.pmoment} with $k = 2$, we have
\begin{equation*}  %\label{e.}
\frac 1 2 \dr_t \E[\td X_t^2] = (1-m_t^2) \E[\td X_t^2] - m_t \E[\td X_t^3] - \frac 1 3 \E[\td X_t^4] + \ep.
\end{equation*}
We recall from \eqref{e.drX2.crude.bound} that $\dr_t \E[\td X_t^2] \le C \al^{\frac 1 2 - k} + \ep$, and a small modification of the argument leading to this bound, starting from the previous display, allows one to bound the absolute value of $\dr_t \E[\td X_t^2]$ by the same expression. For $\ep > 0$ sufficiently small, using again that $T_\al \le 4\pi/\al$, we thus have that
\begin{equation*}  %\label{e.}
\int_{t_0}^{t_0 + T_\al}(1-m_t^2)\Ll( \E[\td X_t^2] - \E[\td X_{t_0}^2]\Rr)\, \d t \le C \al^{-\frac 3 2 - k}.
\end{equation*}
Using \eqref{e.ito.pmoment} with $k = 3$ and similar arguments as those leading to the bound on $|\dr_t \E[\td X_t^2]|$, one can show that $|\dr_t \E[\td X_t^3]| \le C \al^{\frac 1 2 - k}$, and therefore
\begin{equation*}  %\label{e.}
\int_{t_0}^{t_0 + T_\al} m_t \Ll( \E[\td X_{t_0}^3] - \E[\td X_{t}^3] \Rr) \, \d t \le C \al^{-\frac 3 2 - k}.
\end{equation*}
Combining the last three displays (and using again that we take $\ep > 0$ sufficiently small and adjusting $C$), we obtain that 
\begin{multline*}  %\label{e.}
\E[\td X_{t_0 + T_\al}^2] 
\le \E[\td X_{t_0}^2]\Ll(1 + 2 \int_{t_0}^{t_0 + T_\al} (1-m_t^2)\, \d t \Rr) 
\\
- 2 \E[\td X_{t_0}^3] \int_{t_0}^{t_0 + T_\al} m_t \, \d t + C \al^{-\frac 3 2 - k}.
\end{multline*}
Using again that $|\E[\td X_{t_0}^3]| \le C \al^{\frac 1 2 - k}$ and \eqref{e.nice.m} yields that
\begin{equation*}  %\label{e.}
\E[\td X_{t_0 + T_\al}^2] \le \E[\td X_{t_0}^2] \Ll( 1 - \frac{T_\al}{2} \Rr) + C \al^{\frac 1 2 - 2k} + C \al^{-\frac 3 2 - k}.
\end{equation*}
Recalling that $\E[\td X_{t_0}^2] \le \al^{-k}$, that $T_\al \ge 2\pi/\al$, and that $k > \frac 3 2$, we obtain the desired result by taking $\al$ sufficiently large.
\end{proof}

\end{document}
